\documentclass[10pt,a4paper]{amsart}
\usepackage[margin=19mm]{geometry}
\usepackage{amsmath,amssymb,mathtools}
\usepackage[scr=rsfs,cal=euler]{mathalfa}
\usepackage{xfrac}
\usepackage[unicode,hidelinks,bookmarks=false]{hyperref}
\newcommand{\ie}{i.e.\ }
\newcommand{\cf}{cf.\ }
\newcommand{\resp}{respectively\ }
\newcommand{\Subsection}{\subsection}
\providecommand{\nobreakdash}{\nobreak\hbox{-}}
\setlength{\emergencystretch}{2em}
\multlinegap0pt
\usepackage[shortlabels]{enumitem}
\usepackage{amssymb}
\usepackage{xcolor}
\usepackage{algorithm}\usepackage{algpseudocode}
\newtheorem{theorem}{Theorem}
\newtheorem{corollary}{Corollary}
\title{Image transformations, Markov operators, and sample median}
\author{S. V. Butler}
\date{Source: arXiv:2605.03130v1 (CC BY 4.0)}
\begin{document}
\maketitle

\noindent\emph{Source and licence.} Image transformations, Markov operators, and sample median. Version arXiv:2605.03130v1 (CC BY 4.0), distributed by arXiv under the Creative Commons Attribution 4.0 International licence, http://creativecommons.org/licenses/by/4.0/, as recorded in the arXiv licence metadata for this exact version and shown on the abstract page https://arxiv.org/abs/2605.03130v1. Full licence notice: \texttt{licenses/arxiv-2605.03130-CC-BY-4.0.txt}.\par\smallskip

\noindent\textit{Editorial scope.} Counted unit: the article's theorem ITcontin (source0.tex lines 1201--1210). The article's complete original proof of it is delivered with it (source0.tex lines 1212--1247, one \texttt{proof} environment of 36 lines). No mathematics is written, repaired or re-derived: the statement and the proof are the article's own lines, in the article's own notation, and no retained line is substituted or re-flowed.  Retained uncounted as the source-local context the proof needs: lines 739--746; lines 760--793; lines 1160--1166; lines 1190--1192.  Nothing was omitted from any retained span, so the package carries no omission record.  No retained line contains a citation, so the wrapper carries no bibliography and no bibliography entry.  No retained line contains a figure environment, an image-inclusion command or drawing code, so no image is delivered and the package declares no asset dependency.  Editorial formatting only, in addition: an AMS \texttt{amsart} preamble that restates the article's own declarations and macros used by the retained lines, namely the environments \texttt{theorem}, \texttt{corollary} on separate counters, together with the standard packages those lines need.  The retained environments print in this document as Theorem 1, Corollary 1 in order of appearance, whereas the article prints the same items with the numbering of its own full document.  The complete native source of the article as distributed in the arXiv e-print for this exact version is bundled unchanged as \texttt{sources/199774/source0.tex}.
\begin{theorem} \label{adjCont}
Suppose $X,Y$ are LC and $q$ is a d-image transformation (respectively, an image transformation) from $X$ to $Y$.
Then the adjoint map $q^* : (DTM(Y), wk) \rightarrow (DTM(X), wk) $ (resp.,  $q^* : (TM(Y), wk) \rightarrow (TM(X), wk) $) 
preserves finite linear combinations and  
is continuous (i.e. $ \mu_t \Longrightarrow \mu$ implies $q^* \mu_t \Longrightarrow  q^*\mu$).
In particular,  $q^* :Y^{\flat} \rightarrow X^{\flat}$ (resp.,  $q^* :Y^\sharp \rightarrow X^\sharp$)  is continuous.
If $y \in q(\{ x\})$ then $q^*(\delta_y) = \delta_x$.
\end{theorem}

\begin{theorem} \label{ITqh}
Suppose $X$ and $Y$  are LC. There is a 1-1 correspondence between d-image transformations (respectively, image transformations) 
$q$ from $X$ to $Y$ and 
conic quasi-homomorphisms $\theta : C_0^+(X) \longrightarrow C_0^+(Y)$ 
(resp., quasi-homomorphisms $\theta : C_0(X) \longrightarrow C_0(Y)$),
and there is a 1-1 correspondence between d-image transformations (resp., image transformations) 
$q$ and continuous k-proper functions $w: Y \longrightarrow X^{\flat}$ (resp.,   $w: Y \longrightarrow X^\sharp$).
Moreover,
\begin{enumerate}[leftmargin=0.25in, label=(h\arabic*),ref=(h\arabic*)]
%\begin{enumerate}
\item \label{ITqh1}
$\theta(f)(y) = \int_X f\, dw_y =  \int_X f \, d(q^*\delta_y)$ (where $w_y = w(y)$).
\item \label{ITqh6}
$ \| \theta(f) \| \le \| f\|$. 
\item \label{ITqh6a}
If $ f, g \in C_c^+(X)$ then $ \| \theta(f) - \theta(g) \| \le \| f -g\|$, and if $X$ is compact, then 
$ \| \theta(f) - \theta(g) \| \le \| f -g\|$ for any $f,g \in C(X)$. 
For a quasi-homomorphism, $\| \theta(f) - \theta(g) \| \le 2 \| f -g\|$. 
\item \label{ITqh4}
For an image transformation $q$ and an open set $ A \subseteq \mathbb{R}$ or a closed set $A \subseteq \mathbb{R} \setminus \{ 0\}$ 
\begin{align} \label{qthetaf}
q(f^{-1}(A)) = (\theta(f))^{-1} (A),
\end{align}
and for compact $X$ this holds for any open or closed $ A \subseteq \mathbb{R}$. For a d-image transformation $q$ equality (\ref{qthetaf}) holds for
$A=[t, \infty)$ and $A=(t, \infty)$ for  almost every $t$ .
\item \label{ITqh12}
$q(X) = Y$.
\item \label{ITqh15}
$(q(E))^{\flat}  = (q^*)^{-1}(E^{\flat})$ (resp., $(q(E))^{\sharp}  = (q^*)^{-1}(E^{\sharp})$ for a compact or an open set $E \subseteq X$.
\item \label{ITqh5}
%For $f \in C_0(X)$ and a quasi-linear map $\theta$ (resp., for $f \in C_0^+(X)$ and a quasi-homomorphism $ \theta$)
$ \int_X f \, dq^*\nu = \int_Y \theta(f) \, d\nu \  $.
\end{enumerate}
\end{theorem}

\begin{corollary} \label{inftSuDTM}
Suppose $X$ is locally compact, 
family $\mathcal{M}$ is a uniformly bounded in variation family of (deficient) topological measures on $X$, and 
$\sum_{i=1}^\infty \alpha_i \le 1, \, \alpha_i \ge 0$.
Then $\mu = \sum_{i=1}^\infty  \alpha_1 \mu_i$  belongs to $\mathcal{M}$,  and 
$ \int f \, d \mu = \sum_{i=1}^\infty  \alpha_i  \int f\, d \mu_i$ for any $ f \in C_0(X)$.  
\end{corollary}

\begin{align} \label{MarkovS} 
S(\mu) = \sum_{i=1}^\infty \alpha_i q_i^* (\mu). 
\end{align}

\begin{theorem} \label{ITcontin}
Suppose $X, Y$ are locally compact metric spaces,   $\mathcal{M}$ is a uniformly bounded in variation family of (deficient) topological measures on $Y$. 
Let $\{ X, Y, q_i^*, \alpha_i \} $ be a (d-) image transformation system with  $\alpha_i \ge 0$ and $\sum_{i=1}^\infty \alpha_i < \infty$.
Then the operator $S: (\mathcal{M}, wk) \longrightarrow (TM(X), wk)$ (resp.,  $S: (\mathcal{M}, wk) \longrightarrow (DTM(X), wk)$) 
given by (\ref{MarkovS}) 
is a continuous Markov operator which preserves countable convex linear combinations. 
The image of $S$ is a uniformly bounded in variation family of (deficient) topological measures on $X$. 
On the collection of measures from $\mathcal{M}$, $S$ is a Markov-Feller operator with nonlinear
dual operator $T = \sum_{i=1}^\infty \alpha_i \theta_i$, where each $\theta_i$ is a (conic) quasi-homomorphism corresponding to $q_i$. 
\end{theorem}

\begin{proof}
Let  $ \| \mu \|  \le L$ for each $ \mu \in \mathcal{M}$. Then $S(\mu)(X) \le L  \sum_{i=1}^\infty \alpha_i $, 
so $S(\mathcal{M})$ is a uniformly bounded in variation family of (deficient) topological measures on $X$.
By Corollary \ref{inftSuDTM}, for $ f \in C_0(X)$ 
\begin{align} \label{INTdSmu}
 \int_X f \, d(S(\mu)) = \sum_{i=1}^\infty \alpha_i \int_X f \, d(q_i^* \mu).
\end{align}

By Corollary \ref{inftSuDTM},  a countable convex linear combination of members of $\mathcal{M}$ is in $\mathcal{M}$, 
and it is easy to see that $S$ preserves countable convex linear combinations. 
Suppose $ \mu_t \Longrightarrow \mu$. 
Let $ f \in C_0(X)$. For $ \epsilon >0$ pick $n$ such that 
$ L \| f \| \sum_{i=n+1}^{\infty} \alpha_i  < \epsilon$. 
Then $| \sum_{i=n+1}^\infty \alpha_i \int_X f \, d(q_i^* \nu)| < \epsilon$ for any $ \nu \in  \mathcal{M} $.
By Theorem \ref{adjCont} each  $q_i^*$ is continuous, so there is $ t_0$ such that 
$ \sum_{i=1}^n \alpha_i  |  \int_X f \, d(q_i^* \mu_t) -  \int_X f \, d(q_i^* \mu) | < \epsilon$ for $ t \ge t_0$.
Then for  $ t \ge t_0$
\begin{align*}
| \int_X f \, d(S(\mu_t)) -  \int_X f \, d(S(\mu)) | &= 
|\sum_{i=1}^{\infty} \alpha_i \int_X f \, d(q_i^*\mu_t) -  \sum_{i=1}^{\infty} \alpha_i \int_X f \, d(q_i^*\mu) | \\
& \le   \sum_{i=1}^n |  \alpha_i   \int_X f \, d(q_i^* \mu_t) -  \int_X f \, d(q_i^* \mu) | + 2 \epsilon < 3 \epsilon. 
\end{align*}
%\begin{align*}
%| \int_X f\, d(S(\mu_{\alpha})) -  \int_X f\, d(S(\mu))| &= 
%|\sum_{i=1}^{\infty} \alpha_i \int_X f\, d(q^*(\mu_{\alpha}) -  \sum_{i=1}^\infty \alpha_i \int_X f\, d(q^*(\mu)) \\
%&\le \sum_{i=1}^N \alpha_i  |  \int_X f\, d(q^* \mu_{\alpha}) -  \int_X f\, d(q^* \mu) | + | \sum_{i=N+1}^\infty \alpha_i \int_X f \, d(q^* \mu) | + 
%| \sum_{i=N+1}^\infty \alpha_i \int_X f \, d(q^* \mu_{\alpha}) < 3 eps,
%\end{align*}
Thus,  $S( \mu_t) \Longrightarrow  S(\mu)$, and $S$ is continuous. 

Now let $ \mu$ be a measure from $\mathcal{M}$.
From formula (\ref{INTdSmu}) and Theorem \ref{ITqh}\ref{ITqh5} we have:
$$ \int_X f \, d(S(\mu)) = \sum_{i=1}^\infty \alpha_i \int_X f \, d(q_i^* \mu) =  \sum_{i=1}^\infty \int_Y  \alpha_i  \theta_i(f) \,  d\mu 
= \int_Y  \sum_{i=1}^\infty  \alpha_i  \theta_i(f) \,  d\mu,$$
so $ T= \sum_{i=1}^\infty \alpha_i \theta_i$ is the dual operator for $S$. $T$ is nonlinear, since $\theta_i$ are nonlinear. 
\end{proof}

\end{document}
